% Options for packages loaded elsewhere
% Options for packages loaded elsewhere
\PassOptionsToPackage{unicode}{hyperref}
\PassOptionsToPackage{hyphens}{url}
\PassOptionsToPackage{dvipsnames,svgnames,x11names}{xcolor}
%
\documentclass[
  letterpaper,
  DIV=11,
  numbers=noendperiod]{scrartcl}
\usepackage{xcolor}
\usepackage{amsmath,amssymb}
\setcounter{secnumdepth}{5}
\usepackage{iftex}
\ifPDFTeX
  \usepackage[T1]{fontenc}
  \usepackage[utf8]{inputenc}
  \usepackage{textcomp} % provide euro and other symbols
\else % if luatex or xetex
  \usepackage{unicode-math} % this also loads fontspec
  \defaultfontfeatures{Scale=MatchLowercase}
  \defaultfontfeatures[\rmfamily]{Ligatures=TeX,Scale=1}
\fi
\usepackage{lmodern}
\ifPDFTeX\else
  % xetex/luatex font selection
\fi
% Use upquote if available, for straight quotes in verbatim environments
\IfFileExists{upquote.sty}{\usepackage{upquote}}{}
\IfFileExists{microtype.sty}{% use microtype if available
  \usepackage[]{microtype}
  \UseMicrotypeSet[protrusion]{basicmath} % disable protrusion for tt fonts
}{}
\makeatletter
\@ifundefined{KOMAClassName}{% if non-KOMA class
  \IfFileExists{parskip.sty}{%
    \usepackage{parskip}
  }{% else
    \setlength{\parindent}{0pt}
    \setlength{\parskip}{6pt plus 2pt minus 1pt}}
}{% if KOMA class
  \KOMAoptions{parskip=half}}
\makeatother
% Make \paragraph and \subparagraph free-standing
\makeatletter
\ifx\paragraph\undefined\else
  \let\oldparagraph\paragraph
  \renewcommand{\paragraph}{
    \@ifstar
      \xxxParagraphStar
      \xxxParagraphNoStar
  }
  \newcommand{\xxxParagraphStar}[1]{\oldparagraph*{#1}\mbox{}}
  \newcommand{\xxxParagraphNoStar}[1]{\oldparagraph{#1}\mbox{}}
\fi
\ifx\subparagraph\undefined\else
  \let\oldsubparagraph\subparagraph
  \renewcommand{\subparagraph}{
    \@ifstar
      \xxxSubParagraphStar
      \xxxSubParagraphNoStar
  }
  \newcommand{\xxxSubParagraphStar}[1]{\oldsubparagraph*{#1}\mbox{}}
  \newcommand{\xxxSubParagraphNoStar}[1]{\oldsubparagraph{#1}\mbox{}}
\fi
\makeatother

\usepackage{color}
\usepackage{fancyvrb}
\newcommand{\VerbBar}{|}
\newcommand{\VERB}{\Verb[commandchars=\\\{\}]}
\DefineVerbatimEnvironment{Highlighting}{Verbatim}{commandchars=\\\{\}}
% Add ',fontsize=\small' for more characters per line
\usepackage{framed}
\definecolor{shadecolor}{RGB}{241,243,245}
\newenvironment{Shaded}{\begin{snugshade}}{\end{snugshade}}
\newcommand{\AlertTok}[1]{\textcolor[rgb]{0.68,0.00,0.00}{#1}}
\newcommand{\AnnotationTok}[1]{\textcolor[rgb]{0.37,0.37,0.37}{#1}}
\newcommand{\AttributeTok}[1]{\textcolor[rgb]{0.40,0.45,0.13}{#1}}
\newcommand{\BaseNTok}[1]{\textcolor[rgb]{0.68,0.00,0.00}{#1}}
\newcommand{\BuiltInTok}[1]{\textcolor[rgb]{0.00,0.23,0.31}{#1}}
\newcommand{\CharTok}[1]{\textcolor[rgb]{0.13,0.47,0.30}{#1}}
\newcommand{\CommentTok}[1]{\textcolor[rgb]{0.37,0.37,0.37}{#1}}
\newcommand{\CommentVarTok}[1]{\textcolor[rgb]{0.37,0.37,0.37}{\textit{#1}}}
\newcommand{\ConstantTok}[1]{\textcolor[rgb]{0.56,0.35,0.01}{#1}}
\newcommand{\ControlFlowTok}[1]{\textcolor[rgb]{0.00,0.23,0.31}{\textbf{#1}}}
\newcommand{\DataTypeTok}[1]{\textcolor[rgb]{0.68,0.00,0.00}{#1}}
\newcommand{\DecValTok}[1]{\textcolor[rgb]{0.68,0.00,0.00}{#1}}
\newcommand{\DocumentationTok}[1]{\textcolor[rgb]{0.37,0.37,0.37}{\textit{#1}}}
\newcommand{\ErrorTok}[1]{\textcolor[rgb]{0.68,0.00,0.00}{#1}}
\newcommand{\ExtensionTok}[1]{\textcolor[rgb]{0.00,0.23,0.31}{#1}}
\newcommand{\FloatTok}[1]{\textcolor[rgb]{0.68,0.00,0.00}{#1}}
\newcommand{\FunctionTok}[1]{\textcolor[rgb]{0.28,0.35,0.67}{#1}}
\newcommand{\ImportTok}[1]{\textcolor[rgb]{0.00,0.46,0.62}{#1}}
\newcommand{\InformationTok}[1]{\textcolor[rgb]{0.37,0.37,0.37}{#1}}
\newcommand{\KeywordTok}[1]{\textcolor[rgb]{0.00,0.23,0.31}{\textbf{#1}}}
\newcommand{\NormalTok}[1]{\textcolor[rgb]{0.00,0.23,0.31}{#1}}
\newcommand{\OperatorTok}[1]{\textcolor[rgb]{0.37,0.37,0.37}{#1}}
\newcommand{\OtherTok}[1]{\textcolor[rgb]{0.00,0.23,0.31}{#1}}
\newcommand{\PreprocessorTok}[1]{\textcolor[rgb]{0.68,0.00,0.00}{#1}}
\newcommand{\RegionMarkerTok}[1]{\textcolor[rgb]{0.00,0.23,0.31}{#1}}
\newcommand{\SpecialCharTok}[1]{\textcolor[rgb]{0.37,0.37,0.37}{#1}}
\newcommand{\SpecialStringTok}[1]{\textcolor[rgb]{0.13,0.47,0.30}{#1}}
\newcommand{\StringTok}[1]{\textcolor[rgb]{0.13,0.47,0.30}{#1}}
\newcommand{\VariableTok}[1]{\textcolor[rgb]{0.07,0.07,0.07}{#1}}
\newcommand{\VerbatimStringTok}[1]{\textcolor[rgb]{0.13,0.47,0.30}{#1}}
\newcommand{\WarningTok}[1]{\textcolor[rgb]{0.37,0.37,0.37}{\textit{#1}}}

\usepackage{longtable,booktabs,array}
\newcounter{none} % for unnumbered tables
\usepackage{calc} % for calculating minipage widths
% Correct order of tables after \paragraph or \subparagraph
\usepackage{etoolbox}
\makeatletter
\patchcmd\longtable{\par}{\if@noskipsec\mbox{}\fi\par}{}{}
\makeatother
% Allow footnotes in longtable head/foot
\IfFileExists{footnotehyper.sty}{\usepackage{footnotehyper}}{\usepackage{footnote}}
\makesavenoteenv{longtable}
\usepackage{graphicx}
\makeatletter
\newsavebox\pandoc@box
\newcommand*\pandocbounded[1]{% scales image to fit in text height/width
  \sbox\pandoc@box{#1}%
  \Gscale@div\@tempa{\textheight}{\dimexpr\ht\pandoc@box+\dp\pandoc@box\relax}%
  \Gscale@div\@tempb{\linewidth}{\wd\pandoc@box}%
  \ifdim\@tempb\p@<\@tempa\p@\let\@tempa\@tempb\fi% select the smaller of both
  \ifdim\@tempa\p@<\p@\scalebox{\@tempa}{\usebox\pandoc@box}%
  \else\usebox{\pandoc@box}%
  \fi%
}
% Set default figure placement to htbp
\def\fps@figure{htbp}
\makeatother


% definitions for citeproc citations
\NewDocumentCommand\citeproctext{}{}
\NewDocumentCommand\citeproc{mm}{%
  \begingroup\def\citeproctext{#2}\cite{#1}\endgroup}
\makeatletter
 % allow citations to break across lines
 \let\@cite@ofmt\@firstofone
 % avoid brackets around text for \cite:
 \def\@biblabel#1{}
 \def\@cite#1#2{{#1\if@tempswa , #2\fi}}
\makeatother
\newlength{\cslhangindent}
\setlength{\cslhangindent}{1.5em}
\newlength{\csllabelwidth}
\setlength{\csllabelwidth}{3em}
\newenvironment{CSLReferences}[2] % #1 hanging-indent, #2 entry-spacing
 {\begin{list}{}{%
  \setlength{\itemindent}{0pt}
  \setlength{\leftmargin}{0pt}
  \setlength{\parsep}{0pt}
  % turn on hanging indent if param 1 is 1
  \ifodd #1
   \setlength{\leftmargin}{\cslhangindent}
   \setlength{\itemindent}{-1\cslhangindent}
  \fi
  % set entry spacing
  \setlength{\itemsep}{#2\baselineskip}}}
 {\end{list}}
\usepackage{calc}
\newcommand{\CSLBlock}[1]{\hfill\break\parbox[t]{\linewidth}{\strut\ignorespaces#1\strut}}
\newcommand{\CSLLeftMargin}[1]{\parbox[t]{\csllabelwidth}{\strut#1\strut}}
\newcommand{\CSLRightInline}[1]{\parbox[t]{\linewidth - \csllabelwidth}{\strut#1\strut}}
\newcommand{\CSLIndent}[1]{\hspace{\cslhangindent}#1}



\setlength{\emergencystretch}{3em} % prevent overfull lines

\providecommand{\tightlist}{%
  \setlength{\itemsep}{0pt}\setlength{\parskip}{0pt}}



 


\KOMAoption{captions}{tableheading}
\makeatletter
\@ifpackageloaded{tcolorbox}{}{\usepackage[skins,breakable]{tcolorbox}}
\@ifpackageloaded{fontawesome5}{}{\usepackage{fontawesome5}}
\definecolor{quarto-callout-color}{HTML}{909090}
\definecolor{quarto-callout-note-color}{HTML}{0758E5}
\definecolor{quarto-callout-important-color}{HTML}{CC1914}
\definecolor{quarto-callout-warning-color}{HTML}{EB9113}
\definecolor{quarto-callout-tip-color}{HTML}{00A047}
\definecolor{quarto-callout-caution-color}{HTML}{FC5300}
\definecolor{quarto-callout-color-frame}{HTML}{acacac}
\definecolor{quarto-callout-note-color-frame}{HTML}{4582ec}
\definecolor{quarto-callout-important-color-frame}{HTML}{d9534f}
\definecolor{quarto-callout-warning-color-frame}{HTML}{f0ad4e}
\definecolor{quarto-callout-tip-color-frame}{HTML}{02b875}
\definecolor{quarto-callout-caution-color-frame}{HTML}{fd7e14}
\makeatother
\makeatletter
\@ifpackageloaded{caption}{}{\usepackage{caption}}
\AtBeginDocument{%
\ifdefined\contentsname
  \renewcommand*\contentsname{Table of contents}
\else
  \newcommand\contentsname{Table of contents}
\fi
\ifdefined\listfigurename
  \renewcommand*\listfigurename{List of Figures}
\else
  \newcommand\listfigurename{List of Figures}
\fi
\ifdefined\listtablename
  \renewcommand*\listtablename{List of Tables}
\else
  \newcommand\listtablename{List of Tables}
\fi
\ifdefined\figurename
  \renewcommand*\figurename{Figure}
\else
  \newcommand\figurename{Figure}
\fi
\ifdefined\tablename
  \renewcommand*\tablename{Table}
\else
  \newcommand\tablename{Table}
\fi
}
\@ifpackageloaded{float}{}{\usepackage{float}}
\floatstyle{ruled}
\@ifundefined{c@chapter}{\newfloat{codelisting}{h}{lop}}{\newfloat{codelisting}{h}{lop}[chapter]}
\floatname{codelisting}{Listing}
\newcommand*\listoflistings{\listof{codelisting}{List of Listings}}
\makeatother
\makeatletter
\makeatother
\makeatletter
\@ifpackageloaded{caption}{}{\usepackage{caption}}
\@ifpackageloaded{subcaption}{}{\usepackage{subcaption}}
\makeatother
\usepackage{bookmark}
\IfFileExists{xurl.sty}{\usepackage{xurl}}{} % add URL line breaks if available
\urlstyle{same}
\hypersetup{
  pdftitle={Quantile regression of gestational age and preterm birth for epidemiology research},
  pdfauthor={Lawrence McCandless},
  pdfkeywords={Quantile regression, gestational age, preterm
birth, pregnancy outcomes, epidemiology, statistical modeling},
  colorlinks=true,
  linkcolor={blue},
  filecolor={Maroon},
  citecolor={Blue},
  urlcolor={Blue},
  pdfcreator={LaTeX via pandoc}}


\title{Quantile regression of gestational age and preterm birth for
epidemiology research}
\usepackage{etoolbox}
\makeatletter
\providecommand{\subtitle}[1]{% add subtitle to \maketitle
  \apptocmd{\@title}{\par {\large #1 \par}}{}{}
}
\makeatother
\subtitle{Tutorial Essay and Methodological Note}
\author{Lawrence McCandless}
\date{2026-05-01}
\begin{document}
\maketitle
\begin{abstract}
Quantile regression is described for the epidemiology of gestational age
and preterm birth. Unlike ordinary least squares regression, which
models only the mean outcome, quantile regression models the entire
gestational age distribution, including the lower tails associated with
preterm birth. I illustrate the interpretation of quantile regression
results and implementation in the R package \texttt{quantreg}. Using
data from the 1958 UK Perinatal Mortality Survey, maternal albuminuria
is shown to be associated with shorter gestational age, particularly
among pregnancies in the lower quantiles of the gestational age
distribution. Conventional quantile regression is also shown to produce
implausible integer-valued coefficient estimates when gestational age is
analyzed as a discrete outcome measured in days.
\end{abstract}

\renewcommand*\contentsname{Table of contents}
{
\setcounter{tocdepth}{3}
\tableofcontents
}

\begin{tcolorbox}[enhanced jigsaw, arc=.35mm, bottomrule=.15mm, breakable, colback=white, colframe=quarto-callout-note-color-frame, left=2mm, leftrule=.75mm, opacityback=0, rightrule=.15mm, toprule=.15mm]
\begin{minipage}[t]{5.5mm}
\textcolor{quarto-callout-note-color}{\faInfo}
\end{minipage}%
\begin{minipage}[t]{\textwidth - 5.5mm}

\textbf{Lawrence McCandless} is Professor, Faculty of Health Sciences,
Simon Fraser University, 8888 University Drive, Burnaby BC V5A 1S6,
Canada.

E-mail: \texttt{lmccandl@sfu.ca}

R codes and a synthetic version of the PMS dataset used are available
from the GitHub repository for this monograph:

\url{https://github.com/Lawrence-McCandless/QuantileRegressionEpidemiology}

\end{minipage}%
\end{tcolorbox}

\textbf{Key words and phrases:} Quantile regression, gestational age,
preterm birth, pregnancy outcomes, epidemiology, statistical modeling.

\section{Quantile Regression: The Basic
Idea}\label{quantile-regression-the-basic-idea}

Quantile regression is a data analysis technique where quantiles of the
outcome are modelled in relation to predictor variables (Davino et al.
2014; Koenker 2005). It is an extension of least squares regression,
which models averages. Quantile regression gives a richer description of
the distribution of the outcome. It provides information about the
entire distribution of the response, and this allows for a deeper
understanding of how predictor variables influence different parts of
the outcome distribution, not just the average.

This is useful in epidemiology research ({Strickland et al.} 2009). For
example, in child health studies we can model IQ scores for
neurodevelopment, or the forced expiratory volume to measure lung
function. In these settings, we might prefer to not use linear
regression models for location and scale because it assumes a normally
distributed model for the outcome data distribution. Quantile regression
is more robust to heteroscedasticity and outliers.

The basic idea of quantile regression is that the investigator chooses a
particular quantile, which is indexed the Greek letter \(\tau\), which
ranges from zero to one (from 0\% to 100\%). Informally, a quantile is
simply the same thing as a ``percentile'', except that a percentile is a
percentage whereas a quantile is a number between zero and one. If
\(\tau=0.5\) then the quantile is the median. If \(\tau=0.1\), then the
quantile is the 10\(^{th}\) percentile, meaning the bottom 10\% of the
distribution.

Quantile regression involves modeling the quantile of a dependent
variable \(Y\) in relation to a predictor variable \(X\). The model is
written \[\begin{aligned}
Q_{\tau} = \beta_{0 [\tau]} + \beta_{X [\tau]} X, \label{eq-fqr-qrmodel.univar.with.taus}
\end{aligned}\] which has two unknown parameters that define the model:
\(\beta_{0 [\tau]}\) and \(\beta_{X [\tau]}\). The quantity
\(\beta_{X [\tau]}\) is the slope (regression coefficient) and
\(\beta_{0 [\tau]}\) is the y-intercept. The subscript \(\tau\) is
written to emphasize that the parameter depends on \(\tau\). For
example, if \(\tau=0.1\), then slope \(\beta_{X [\tau=0.1]}\) describes
the change in the 0.1 quantile (the 10\(^{th}\) percentile) for every
one-unit increase in \(X\). The y-intercept parameter
\(\beta_{0 [\tau]}\) is equal to the \(\tau^{th}\) quantile in the group
with \(X=0\).

The true values of \(\beta_{0 [\tau]}\) and \(\beta_{X [\tau]}\) are
unknown, and they must be estimated from the data
\(Y_1, Y_2, \dots, Y_n\) and \(X_1, X_2, \dots, X_n\). The conventional
approach described here is to minimize the sum of the weighted absolute
value of the residuals. This is called minimizing the quantile loss, or
equivalently, the pinball loss because the absolute value of the
residuals looks like a pinball machine. To explain in more detail,
quantile regression calculates the quantities \(\hat{\beta}_{0 [\tau]}\)
and \(\hat{\beta}_{X [\tau]}\) which make the sum
\begin{equation}\protect\phantomsection\label{eq-fqr-QRobjFunction}{
\begin{aligned}
\sum_{i=1}^n \rho_\tau\Big( Y_i - (\hat{\beta}_{0 [\tau]} + \hat{\beta}_{X [\tau]} X_i) \Big)
\end{aligned}
}\end{equation} as small as possible, where \(\rho_\tau(.)\) is the
pinball loss function, and it is the absolute value of the residual
\(Y_i - (\beta_{0 [\tau]} + \beta_{X [\tau]} X_i)\) multiplied by
\(\tau\) or \(1-\tau\).

\[
\rho_\tau(u) =
\begin{cases}
\tau |u|, & \text{if } u \geq 0, \\
(1-\tau)|u|, & \text{if } u < 0.
\end{cases}
\]

The minimization in equation Equation~\ref{eq-fqr-QRobjFunction} is
called the quantile regression loss function, and we choose the model
that minimizes the loss. This results in estimating quantiles. To
understand how the model fitting works, it is useful to compare quantile
regression to ordinary least squares regression where we minimize the
sum of the squares of the residuals, which is written
\[\sum_{i=1}^n \Big( Y_i - (\hat{\beta}_{0 [\tau]} + \hat{\beta}_{X [\tau]} X_i) \Big)^2.\]
We illustrate quantile regression using the R package \texttt{quantreg}.

\section{Example: Gestational age and preterm birth in the 1958
Perinatal Mortality Survey
(PMS)}\label{example-gestational-age-and-preterm-birth-in-the-1958-perinatal-mortality-survey-pms}

Early gestational age is an important determinant of infant morbidity
and mortality (Goldenberg et al. 2008). It is measured on a discrete
scale where the unit of analysis is typically weeks or days. Gestational
age is not symmetrically distributed, and it has a long left tail, which
corresponds to infants born too soon. At the upper end of the
distribution, the right tail ends abruptly due to the modern obstetric
practice of induction of labour late in pregnancy. In perinatal
epidemiology, we are interested in the entire shape of the distribution,
not just the mean.

We illustrate these ideas using data from the 1958 UK Perinatal
Mortality Study (PMS) from the UK Data Service (Centre for Longitudinal
Studies 2024). This was a longitudinal cohort study following the lives
of several thousand individuals born in a single week in England,
Scotland and Wales. The data are contained in the file \texttt{pms.csv},
which is available from the GitHub repository of this article. The data
include information from 2,540 mother and child pairs, and measurements
of gestational age, measured in days, are contained in the
\texttt{gestage} variable. For illustration of quantile regression, we
focus on the influence of albuminuria on gestational age. Albuminuria
(\texttt{album} variable coded as 1 or 0) is a symptom of
pregnancy-induced high blood pressure (preeclampsia), which is a known
risk factor for preterm birth and growth restriction. R codes and a
synthetic version of the PMS dataset are available from the GitHub
repository for this monograph at
\url{https://github.com/Lawrence-McCandless/QuantileRegressionEpidemiology}.

We begin by calculating descriptive statistics for gestational age,
stratified by albuminuria level of the mother during pregnancy. Using
the \texttt{quantile()} command in R, we have

\begin{Shaded}
\begin{Highlighting}[]
\SpecialCharTok{\textgreater{}} \FunctionTok{quantile}\NormalTok{(pms}\SpecialCharTok{$}\NormalTok{ga, }\AttributeTok{probs=}\FunctionTok{c}\NormalTok{(}\FloatTok{0.1}\NormalTok{, }\FloatTok{0.25}\NormalTok{, }\FloatTok{0.5}\NormalTok{, }\FloatTok{0.75}\NormalTok{, }\FloatTok{0.9}\NormalTok{), }\AttributeTok{type=}\DecValTok{1}\NormalTok{)}

\DecValTok{10}\SpecialCharTok{\% 25\%} \DecValTok{50}\SpecialCharTok{\% 75\%} \DecValTok{90}\NormalTok{\%}
\DecValTok{268} \DecValTok{275} \DecValTok{282} \DecValTok{289} \DecValTok{294}
\end{Highlighting}
\end{Shaded}

and the results are summarized in Table
Table~\ref{tbl-fqr-pmsalbumvsga}, which also includes the frequency of
preterm births, defined as births shorter than 259 days (37 weeks).

\begin{longtable}[]{@{}
  >{\raggedright\arraybackslash}p{(\linewidth - 20\tabcolsep) * \real{0.0698}}
  >{\raggedleft\arraybackslash}p{(\linewidth - 20\tabcolsep) * \real{0.0930}}
  >{\raggedleft\arraybackslash}p{(\linewidth - 20\tabcolsep) * \real{0.0930}}
  >{\raggedleft\arraybackslash}p{(\linewidth - 20\tabcolsep) * \real{0.0930}}
  >{\raggedleft\arraybackslash}p{(\linewidth - 20\tabcolsep) * \real{0.0930}}
  >{\raggedleft\arraybackslash}p{(\linewidth - 20\tabcolsep) * \real{0.0930}}
  >{\raggedleft\arraybackslash}p{(\linewidth - 20\tabcolsep) * \real{0.0930}}
  >{\raggedleft\arraybackslash}p{(\linewidth - 20\tabcolsep) * \real{0.0930}}
  >{\raggedleft\arraybackslash}p{(\linewidth - 20\tabcolsep) * \real{0.0930}}
  >{\raggedleft\arraybackslash}p{(\linewidth - 20\tabcolsep) * \real{0.0930}}
  >{\raggedleft\arraybackslash}p{(\linewidth - 20\tabcolsep) * \real{0.0930}}@{}}
\caption{Distribution of gestational age and the frequency of preterm
birth in the PMS dataset, stratified by albuminuria
level.}\label{tbl-fqr-pmsalbumvsga}\tabularnewline
\toprule\noalign{}
\begin{minipage}[b]{\linewidth}\raggedright
Albuminuria
\end{minipage} & \begin{minipage}[b]{\linewidth}\raggedleft
n
\end{minipage} & \begin{minipage}[b]{\linewidth}\raggedleft
Mean
\end{minipage} & \begin{minipage}[b]{\linewidth}\raggedleft
SD
\end{minipage} & \begin{minipage}[b]{\linewidth}\raggedleft
0.1
\end{minipage} & \begin{minipage}[b]{\linewidth}\raggedleft
0.25
\end{minipage} & \begin{minipage}[b]{\linewidth}\raggedleft
0.5
\end{minipage} & \begin{minipage}[b]{\linewidth}\raggedleft
0.75
\end{minipage} & \begin{minipage}[b]{\linewidth}\raggedleft
0.9
\end{minipage} & \begin{minipage}[b]{\linewidth}\raggedleft
count
\end{minipage} & \begin{minipage}[b]{\linewidth}\raggedleft
\%
\end{minipage} \\
\midrule\noalign{}
\endfirsthead
\toprule\noalign{}
\begin{minipage}[b]{\linewidth}\raggedright
Albuminuria
\end{minipage} & \begin{minipage}[b]{\linewidth}\raggedleft
n
\end{minipage} & \begin{minipage}[b]{\linewidth}\raggedleft
Mean
\end{minipage} & \begin{minipage}[b]{\linewidth}\raggedleft
SD
\end{minipage} & \begin{minipage}[b]{\linewidth}\raggedleft
0.1
\end{minipage} & \begin{minipage}[b]{\linewidth}\raggedleft
0.25
\end{minipage} & \begin{minipage}[b]{\linewidth}\raggedleft
0.5
\end{minipage} & \begin{minipage}[b]{\linewidth}\raggedleft
0.75
\end{minipage} & \begin{minipage}[b]{\linewidth}\raggedleft
0.9
\end{minipage} & \begin{minipage}[b]{\linewidth}\raggedleft
count
\end{minipage} & \begin{minipage}[b]{\linewidth}\raggedleft
\%
\end{minipage} \\
\midrule\noalign{}
\endhead
\bottomrule\noalign{}
\endlastfoot
Yes & 182 & 278.3 & 13.2 & 260 & 272 & 280 & 285 & 293 & 10 & 5.5\% \\
No & 2358 & 281.0 & 12.7 & 268 & 275 & 282 & 289 & 295 & 96 & 4.1\% \\
All & 2540 & 280.9 & 12.8 & 268 & 275 & 282 & 289 & 294 & 106 & 4.2\% \\
\end{longtable}

From the results it appears that albuminuria is associated with a
shorter gestational age. Furthermore, the reduction is greatest at the
lower end of the distribution of gestational age. For example, at the
median, we see a 2-day reduction in gestational age (280 days versus 282
days), whereas at the \(\tau=0.1\) quantile (10\(^{th}\) percentile) we
see an 8-day reduction (260 days versus 268 days). A \(>1\) week
reduction in gestational age is clinically significant because the
normal duration of pregnancy is 40 weeks, whereas an infant born at less
than 37 weeks is classified as preterm. Moreover, the shift in the
distribution of gestational age is not apparent based on a comparison of
means. The mean is 278.3 days for mother with albuminuria versus 281.0
days in mothers without. In Table Table~\ref{tbl-fqr-pmsalbumvsga}, we
see that there are a total of 106 preterm births, and that albuminuria
is associated with a 31\% increase in the frequency of preterm birth
(5.5\% versus 4.2\%).

\section{Quantile regression with a single predictor
variable}\label{quantile-regression-with-a-single-predictor-variable}

To motivate quantile regression, Figure~\ref{fig-fqr-pmdescriptive-ga}
shows the gestational age distribution, including the quantile function.
As expected the distribution of gestational age has a long left tail.

\begin{figure}

\centering{

\includegraphics[width=1\linewidth,height=\textheight,keepaspectratio]{index_files/mediabag/figures/fig-fqr-pmdescriptive-ga.pdf}

}

\caption{\label{fig-fqr-pmdescriptive-ga}Distribution of gestational age
in the PMS dataset. Illustration of the probability density function
\(f(x)\), the cumulative distribution function \(F(x)\), and the
quantile function \(F^{-1}(\tau)\). The dashed line indicates the
threshold for preterm birth (\(<259\) days).}

\end{figure}%

To examine the influence of albuminuria on gestational age,
Figure~\ref{fig-fqr-pmsdescriptiveAlbumvsga} shows the distribution of
gestational age stratified by albuminuria level. The red curve is
shifted to the left, which illustrates that infants of mothers with
albuminuria have shorter gestational age.

\begin{figure}

\centering{

\includegraphics[width=1\linewidth,height=\textheight,keepaspectratio]{index_files/mediabag/figures/fig-fqr-pmsdescriptiveAlbumvsga.pdf}

}

\caption{\label{fig-fqr-pmsdescriptiveAlbumvsga}Distribution of
gestational age stratified by albuminuria level. The red curve describes
mothers with albuminuria, whereas the black curve represents mothers
without albuminuria. The dashed line indicates the threshold for preterm
birth (\(<259\) days).}

\end{figure}%

Accordingly, we applied quantile regression to the data with \texttt{ga}
as the dependent variable and \texttt{album} as the predictor variable.
Subsequently, we considered adjustment for additional risk factors for
preterm birth. We use the R function \texttt{rq()} in the
\texttt{quantreg} package to fit quantile regression models.

\begin{Shaded}
\begin{Highlighting}[]
\SpecialCharTok{\textgreater{}} \FunctionTok{summary}\NormalTok{(}\FunctionTok{rq}\NormalTok{(ga }\SpecialCharTok{\textasciitilde{}}\NormalTok{ album, }\AttributeTok{tau=}\FloatTok{0.1}\NormalTok{, }\AttributeTok{data=}\NormalTok{pms))}

\NormalTok{Call}\SpecialCharTok{:} \FunctionTok{rq}\NormalTok{(}\AttributeTok{formula =}\NormalTok{ ga }\SpecialCharTok{\textasciitilde{}}\NormalTok{ album, }\AttributeTok{tau =} \FloatTok{0.1}\NormalTok{, }\AttributeTok{data =}\NormalTok{ pms)}

\NormalTok{tau}\SpecialCharTok{:}\NormalTok{ [}\DecValTok{1}\NormalTok{] }\FloatTok{0.1}

\NormalTok{Coefficients}\SpecialCharTok{:}
\NormalTok{            Value     Std. Error t value   }\FunctionTok{Pr}\NormalTok{(}\SpecialCharTok{\textgreater{}}\ErrorTok{|}\NormalTok{t}\SpecialCharTok{|}\NormalTok{)}
\NormalTok{(Intercept) }\FloatTok{268.00000}   \FloatTok{0.48725}  \FloatTok{550.02221}   \FloatTok{0.00000}
\NormalTok{album        }\SpecialCharTok{{-}}\FloatTok{8.00000}   \FloatTok{2.24580}   \SpecialCharTok{{-}}\FloatTok{3.56221}   \FloatTok{0.00037}
\end{Highlighting}
\end{Shaded}

From the R output, we have the point estimates
\(\hat{\beta}_{X [\tau]}=-8.0\) and \(\hat{\beta}_{0 [\tau]}=268\), and
the equation for the \(\tau=0.1\) quantile is given by
\(Q_{\tau}(Y) = 268.0 - 8.0 X\). By default, the \texttt{rq(.)} command
uses the \texttt{"nid"} option (Davino et al. 2014; Koenker 2005).

The quantile regression results are illustrated in further detail in
Table Table~\ref{tbl-fqr-pmsbetasqr}. For comparison, two other analyses
are also included: 1) ordinary least squares (OLS) regression of
\texttt{ga} on \texttt{album}, which models the mean of gestational age,
and 2) logistic regression analysis of \texttt{preterm} on
\texttt{album}, where preterm birth is defined as gestational age less
than 259 days (\texttt{preterm\ =\ ga\ \textless{}\ 259}).

\begin{longtable}[]{@{}
  >{\raggedright\arraybackslash}p{(\linewidth - 14\tabcolsep) * \real{0.1250}}
  >{\raggedright\arraybackslash}p{(\linewidth - 14\tabcolsep) * \real{0.1250}}
  >{\raggedright\arraybackslash}p{(\linewidth - 14\tabcolsep) * \real{0.1250}}
  >{\raggedright\arraybackslash}p{(\linewidth - 14\tabcolsep) * \real{0.1250}}
  >{\raggedright\arraybackslash}p{(\linewidth - 14\tabcolsep) * \real{0.1250}}
  >{\raggedright\arraybackslash}p{(\linewidth - 14\tabcolsep) * \real{0.1250}}
  >{\raggedright\arraybackslash}p{(\linewidth - 14\tabcolsep) * \real{0.1250}}
  >{\raggedright\arraybackslash}p{(\linewidth - 14\tabcolsep) * \real{0.1250}}@{}}
\caption{Quantile regression, OLS regression, and logistic regression
results. Quantile regression uses \protect\texttt{ga} as the dependent
variable and \protect\texttt{album} as the predictor. The logistic model
uses \protect\texttt{preterm} as the dependent variable. For quantile
regression, the \protect\texttt{rq(.)} command uses the default option
\protect\texttt{se="nid"} to compute standard errors and confidence
intervals.}\label{tbl-fqr-pmsbetasqr}\tabularnewline
\toprule\noalign{}
\begin{minipage}[b]{\linewidth}\raggedright
Method
\end{minipage} & \begin{minipage}[b]{\linewidth}\raggedright
\(\tau\)
\end{minipage} & \begin{minipage}[b]{\linewidth}\raggedright
\(\hat{\beta}_{X[\tau]}\)
\end{minipage} & \begin{minipage}[b]{\linewidth}\raggedright
SE
\end{minipage} & \begin{minipage}[b]{\linewidth}\raggedright
95\% CI
\end{minipage} & \begin{minipage}[b]{\linewidth}\raggedright
\(\hat{\beta}_{0[\tau]}\)
\end{minipage} & \begin{minipage}[b]{\linewidth}\raggedright
SE
\end{minipage} & \begin{minipage}[b]{\linewidth}\raggedright
95\% CI
\end{minipage} \\
\midrule\noalign{}
\endfirsthead
\toprule\noalign{}
\begin{minipage}[b]{\linewidth}\raggedright
Method
\end{minipage} & \begin{minipage}[b]{\linewidth}\raggedright
\(\tau\)
\end{minipage} & \begin{minipage}[b]{\linewidth}\raggedright
\(\hat{\beta}_{X[\tau]}\)
\end{minipage} & \begin{minipage}[b]{\linewidth}\raggedright
SE
\end{minipage} & \begin{minipage}[b]{\linewidth}\raggedright
95\% CI
\end{minipage} & \begin{minipage}[b]{\linewidth}\raggedright
\(\hat{\beta}_{0[\tau]}\)
\end{minipage} & \begin{minipage}[b]{\linewidth}\raggedright
SE
\end{minipage} & \begin{minipage}[b]{\linewidth}\raggedright
95\% CI
\end{minipage} \\
\midrule\noalign{}
\endhead
\bottomrule\noalign{}
\endlastfoot
Quantile regression & 0.9 & -2.00 & 0.95 & (-3.86, -0.14) & 295.00 &
0.37 & (294.28, 295.72) \\
& 0.75 & -4.00 & 1.04 & (-6.04, -1.96) & 289.00 & 0.36 & (288.29,
289.71) \\
& 0.5 & -2.00 & 1.08 & (-4.12, 0.12) & 282.00 & 0.29 & (281.43,
282.57) \\
& 0.25 & -3.00 & 1.33 & (-5.61, -0.39) & 275.00 & 0.27 & (274.47,
275.53) \\
& 0.1 & -8.00 & 2.25 & (-12.40, -3.60) & 268.00 & 0.49 & (267.04,
268.96) \\
OLS regression & - & -2.76 & 0.98 & (-4.69, -0.84) & 281.05 & 0.26 &
(280.53, 281.56) \\
Logistic regression & - & 0.58 & 0.30 & (0.00, 1.16) & -3.07 & 0.10 &
(-3.26, -2.87) \\
\end{longtable}

The results show an overall pattern where albuminuria is associated with
shorter gestational age, particularly among pregnancies in the lower
quantiles of the gestational age distribution. In Table
Table~\ref{tbl-fqr-pmsbetasqr}, the regression coefficients
\(\hat{\beta}_{X [\tau]}\), for various \(\tau\), are both negative and
precise. Moreover, we see the advantages of quantile regression compared
to OLS regression and logistic regression. OLS identifies a
statistically significant inverse association between albuminuria and
gestational age. But the regression coefficient for the difference in
the means (\(\hat{\beta}_X=-2.76\) days) gives a poor description of the
overall relationship in the data. On the other hand, logistic regression
dichotomizes the gestational age variable. This has the advantage of
giving results that are easier to interpret at the expense of discarding
information about the rest of the distribution. The logistic regression
coefficient is 0.58, which corresponds to an odds ratio of 1.79 for the
association between albuminuria and higher odds of preterm birth.
However, the estimate is relatively imprecise with a wide 95\%
confidence interval 1.0 to 3.19.

\section{Implausible results: Illustrating the influence of discrete
outcomes}\label{implausible-results-illustrating-the-influence-of-discrete-outcomes}

A curious feature of Table Table~\ref{tbl-fqr-pmsbetasqr} is that
quantile regression coefficients are always integers (i.e.~numbers like
\(-2,-1, 0, 1, 2, \dots\) etc). This makes the results look implausible
(Alampi et al. 2025). For example, when \(\tau=0.1\) then
\(\hat{\beta}_{X [\tau]} = -8.00\). These integer-valued coefficient
estimates from quantile regress are not an accident. They occurs because
the \texttt{quantreg} software uses the simplex algorithm (Koenker
2005). When \(Y\) and \(X\) are discrete, then the software will only
produce integer valued parameters. The software that minimizes the
pinball loss in equation Equation~\ref{eq-fqr-QRobjFunction} cannot
explore a continuous range of possible values for the unknown parameters
\(\beta_{X [\tau]}\) and \(\beta_{0 [\tau]}\). Instead it considers a
restricted collection of fitted models called basic solutions, which are
the linear models that directly through pairs of data points. The
simplex algorithm works well for quantile regression with continuous
outcomes because the gaps (spacings) between neighboring data points
shrink to zero in large datasets. But it can generate implausible
results for discrete outcomes because the basic solutions are heavily
influenced by discretization arising from rounding of the data. If the
data are discrete then the output of quantile regression software is
also discrete. See Koenker (Koenker 2005), Section 2.2.1 for discussion.

To show why the results are implausible, it is helpful to consider
multivariable quantile regression analysis with several predictor
variables. We include three additional predictors of gestational age:
Child sex, marital status, and smoking status. These are encoded in the
zero-one variables \texttt{female}, \texttt{single}, and zero-one
indicator variables for smoking (\texttt{smoke1} for former smoker
versus non-smoker, and \texttt{smoke2} for smoker versus non-smoker). We
fit the quantile regression model \[\begin{aligned}
Q_{\tau} = \beta_{0 [\tau]} + \beta_{1 [\tau]} album + \beta_{2 [\tau]} female + \beta_{3 [\tau]} single +\beta_{4 [\tau]} smoke1 + \beta_{5 [\tau]} smoke2\label{eq-fqr-qrmodel.mlr}
\end{aligned}\] using the R code:

\begin{Shaded}
\begin{Highlighting}[]
\SpecialCharTok{\textgreater{}} \FunctionTok{summary}\NormalTok{(}\FunctionTok{rq}\NormalTok{(ga }\SpecialCharTok{\textasciitilde{}}\NormalTok{ album }\SpecialCharTok{+}\NormalTok{ female }\SpecialCharTok{+}\NormalTok{ single }\SpecialCharTok{+} \FunctionTok{as.factor}\NormalTok{(smk), }\AttributeTok{tau=}\FloatTok{0.1}\NormalTok{, }\AttributeTok{data=}\NormalTok{pms))}

\NormalTok{Call}\SpecialCharTok{:} \FunctionTok{rq}\NormalTok{(}\AttributeTok{formula =}\NormalTok{ ga }\SpecialCharTok{\textasciitilde{}}\NormalTok{ album }\SpecialCharTok{+}\NormalTok{ female }\SpecialCharTok{+}\NormalTok{ single }\SpecialCharTok{+} \FunctionTok{as.factor}\NormalTok{(smk), }\AttributeTok{tau =} \FloatTok{0.1}\NormalTok{,}
    \AttributeTok{data =}\NormalTok{ pms)}

\NormalTok{tau}\SpecialCharTok{:}\NormalTok{ [}\DecValTok{1}\NormalTok{] }\FloatTok{0.1}

\NormalTok{Coefficients}\SpecialCharTok{:}
\NormalTok{                Value     Std. Error t value   }\FunctionTok{Pr}\NormalTok{(}\SpecialCharTok{\textgreater{}}\ErrorTok{|}\NormalTok{t}\SpecialCharTok{|}\NormalTok{)}
\NormalTok{(Intercept)     }\FloatTok{268.00000}   \FloatTok{0.77273}  \FloatTok{346.82260}   \FloatTok{0.00000}
\NormalTok{album            }\SpecialCharTok{{-}}\FloatTok{6.00000}   \FloatTok{1.88352}   \SpecialCharTok{{-}}\FloatTok{3.18553}   \FloatTok{0.00146}
\NormalTok{female            }\FloatTok{1.00000}   \FloatTok{0.94601}    \FloatTok{1.05707}   \FloatTok{0.29058}
\NormalTok{single           }\SpecialCharTok{{-}}\FloatTok{4.00000}   \FloatTok{4.98359}   \SpecialCharTok{{-}}\FloatTok{0.80263}   \FloatTok{0.42226}
\FunctionTok{as.factor}\NormalTok{(smk)}\DecValTok{1}   \FloatTok{2.00000}   \FloatTok{1.07956}    \FloatTok{1.85261}   \FloatTok{0.06406}
\FunctionTok{as.factor}\NormalTok{(smk)}\DecValTok{2}  \SpecialCharTok{{-}}\FloatTok{2.00000}   \FloatTok{1.20993}   \SpecialCharTok{{-}}\FloatTok{1.65298}   \FloatTok{0.09846}
\end{Highlighting}
\end{Shaded}

Again the regression coefficients are integers. To show this behavior
more clearly, Figure Figure~\ref{fig-fqr-pmsbetasqr.mlr} plots the
estimated quantile regression coefficients for the different variables
when \(\tau\) ranges from zero to one. Random sampling errors are
ignored (e.g.~95\% confidence intervals). We see that the curves bounce
around half-haphazardly because they are stuck on integer values.
Because \(Y\) and \(X\) are discrete, the \texttt{quantreg} package will
only produce integer valued estimates of the model parameters.

\begin{figure}

\centering{

\includegraphics[width=1\linewidth,height=\textheight,keepaspectratio]{index_files/mediabag/figures/fig-fqr-pmsbetasqr.mlr.pdf}

}

\caption{\label{fig-fqr-pmsbetasqr.mlr}Results of multivariable quantile
regression with \ttt{ga} as the dependent variable and \ttt{album},
\ttt{female} and \ttt{smk} as predictor variables. The curves are
quantile regression coefficients where the horizontal axis shows the
different quantiles that range from zero to one.}

\end{figure}%

Consider the gestational age of female infants. It is already known that
females have slightly shorter average duration of pregnancy than males
among full-term births in many human populations (Divon et al. 2002).
But the difference between males and females is small (e.g.~280.6 days
for males versus 279.8 days for females by one estimate (Divon et al.
2002)). Thus we already know that the effect of sex on gestational age
is non-zero. But it is also unlikely to be exactly equal to 1.0 days
because a single day is an arbitrary unit of measurement. In this
example, the size of the target parameters is roughly the same as the
size of the rounding error in the gestational age scale. The simplex
algorithm for fitting quantile regression gives results where the
accuracy of estimates is influenced discreteness of the outcome in way
that is seems unreasonable.

A further consequence of the implausible point estimates is that it
propagates into the calculation of confidence intervals and p-values for
uncertainty assessment. The standard errors by default are produced
using the option \texttt{se="nid"} computed under the assumption that
gestational age is continuous and that the errors are not necessarily
independent and identically distributed. A discussion of statistical
inference in quantile regression is beyond the scope of this monograph,
but the point is the default t-values and p-values from the
\texttt{quantreg} package depend on the implausible point estimates.

\section{Dither: A remedy for implausible results in quantile regression
with discrete
outcomes}\label{dither-a-remedy-for-implausible-results-in-quantile-regression-with-discrete-outcomes}

There are different strategies to fix the implausible results from
quantile regression with discrete outcomes (Alampi et al. 2025). One
option is to use dithering (Machado and Silva 2005). Dithering is a
procedure that adds uniform noise to the outcome variable \(Y\) to
account for the discreteness. For example, if the duration of gestation
is \texttt{ga=}280 days, then we randomly add a number ranging from -0.5
to 0.5 days to 280. The resulting dithered outcome variable is fed into
quantile regression.

To dither the gestational age variable, we use the \texttt{dither()}
command in the \texttt{quantreg} package to create a new variable
\texttt{ga\_d}.

\begin{Shaded}
\begin{Highlighting}[]
\SpecialCharTok{\textgreater{}}\NormalTok{ pms}\SpecialCharTok{$}\NormalTok{ga\_d }\OtherTok{\textless{}{-}} \FunctionTok{dither}\NormalTok{(pms}\SpecialCharTok{$}\NormalTok{ga)}
\SpecialCharTok{\textgreater{}} \FunctionTok{head}\NormalTok{(pms}\SpecialCharTok{$}\NormalTok{ga)}
\NormalTok{[}\DecValTok{1}\NormalTok{] }\DecValTok{303} \DecValTok{261} \DecValTok{285} \DecValTok{298} \DecValTok{277} \DecValTok{303}
\SpecialCharTok{\textgreater{}} \FunctionTok{head}\NormalTok{(pms}\SpecialCharTok{$}\NormalTok{ga\_d)}
\NormalTok{[}\DecValTok{1}\NormalTok{] }\FloatTok{303.2863} \FloatTok{260.9340} \FloatTok{284.9133} \FloatTok{298.4330} \FloatTok{277.4567} \FloatTok{303.3706}
\end{Highlighting}
\end{Shaded}

We then feed the dithered outcome \texttt{ga\_d} into quantile
regression in place of the original variable \texttt{ga}. The results
are given as

\begin{Shaded}
\begin{Highlighting}[]
\SpecialCharTok{\textgreater{}} \FunctionTok{summary}\NormalTok{(}\FunctionTok{rq}\NormalTok{(ga\_d }\SpecialCharTok{\textasciitilde{}}\NormalTok{ album }\SpecialCharTok{+}\NormalTok{ female }\SpecialCharTok{+}\NormalTok{ single }\SpecialCharTok{+} \FunctionTok{as.factor}\NormalTok{(smk), }\AttributeTok{tau=}\FloatTok{0.1}\NormalTok{, }\AttributeTok{data=}\NormalTok{pms))}

\NormalTok{Call}\SpecialCharTok{:} \FunctionTok{rq}\NormalTok{(}\AttributeTok{formula =}\NormalTok{ ga\_d }\SpecialCharTok{\textasciitilde{}}\NormalTok{ album }\SpecialCharTok{+}\NormalTok{ female }\SpecialCharTok{+}\NormalTok{ single }\SpecialCharTok{+} \FunctionTok{as.factor}\NormalTok{(smk),}
    \AttributeTok{tau =} \FloatTok{0.1}\NormalTok{, }\AttributeTok{data =}\NormalTok{ pms)}

\NormalTok{tau}\SpecialCharTok{:}\NormalTok{ [}\DecValTok{1}\NormalTok{] }\FloatTok{0.1}

\NormalTok{Coefficients}\SpecialCharTok{:}
\NormalTok{                Value     Std. Error t value   }\FunctionTok{Pr}\NormalTok{(}\SpecialCharTok{\textgreater{}}\ErrorTok{|}\NormalTok{t}\SpecialCharTok{|}\NormalTok{)}
\NormalTok{(Intercept)     }\FloatTok{267.96566}   \FloatTok{0.79485}  \FloatTok{337.12761}   \FloatTok{0.00000}
\NormalTok{album            }\SpecialCharTok{{-}}\FloatTok{6.51431}   \FloatTok{1.69716}   \SpecialCharTok{{-}}\FloatTok{3.83835}   \FloatTok{0.00013}
\NormalTok{female            }\FloatTok{1.33697}   \FloatTok{1.00523}    \FloatTok{1.33001}   \FloatTok{0.18364}
\NormalTok{single           }\SpecialCharTok{{-}}\FloatTok{3.20120}   \FloatTok{4.76561}   \SpecialCharTok{{-}}\FloatTok{0.67173}   \FloatTok{0.50182}
\FunctionTok{as.factor}\NormalTok{(smk)}\DecValTok{1}   \FloatTok{1.77985}   \FloatTok{1.14689}    \FloatTok{1.55189}   \FloatTok{0.12081}
\FunctionTok{as.factor}\NormalTok{(smk)}\DecValTok{2}  \SpecialCharTok{{-}}\FloatTok{1.87375}   \FloatTok{1.28035}   \SpecialCharTok{{-}}\FloatTok{1.46347}   \FloatTok{0.14346}
\end{Highlighting}
\end{Shaded}

Based on this R output, we can see that the quantile regression results
are more reasonable because the point estimates
\(\hat{\beta}_{1 [\tau]}, \hat{\beta}_{2 [\tau]} \dots , \hat{\beta}_{5 [\tau]}\)
are no longer constrained to lie on the set of integers. Because the
results depend on the specific dithered values of gestational age, it is
recommended to average over the dithered samples (Machado and Silva
2005). This is achieved by using Monte Carlo simulation to replace the
single dithered values with several possible values, and then repeat the
analysis each time to average over uncertainty. The R codes for
dithering are given in the supplementary GitHub repository for this
monograph.

Figure Figure~\ref{fig-fqr-pmsbetasqr.mlr.dithered} shows the results of
multivariable quantile regression after averaging over the dithered
values of gestational age. Thus Figure
Figure~\ref{fig-fqr-pmsbetasqr.mlr.dithered} is identical to Figure
Figure~\ref{fig-fqr-pmsbetasqr.mlr}, except that the results account for
discreteness of gestational age due to rounding. The quantile
coefficient curves are much smoother. Thus the results are more
plausible because the estimated values of
\(\hat{\beta}_{1 [\tau]}, \hat{\beta}_{2 [\tau]} \dots , \hat{\beta}_{5 [\tau]}\)
are not stuck on the set of integers (i.e.~numbers like
\(-2,-1, 0, 1, 2, \dots\)).

\begin{figure}

\centering{

\includegraphics[width=1\linewidth,height=\textheight,keepaspectratio]{index_files/mediabag/figures/fig-fqr-pmsbetasqr.mlr.dithered.pdf}

}

\caption{\label{fig-fqr-pmsbetasqr.mlr.dithered}Quantile regression with
dithering of \ttt{ga}. The curves are quantile regression coefficients
for the predictor variables \ttt{album}, \ttt{female} and \ttt{smk},
where the horizontal axis shows the different quantiles that range from
zero to one.}

\end{figure}%

Returning to the example of the influence of female sex on gestational
age, we see in Figure Figure~\ref{fig-fqr-pmsbetasqr.mlr.dithered} that
the quantile coefficients curve is greater than zero. In contrast,
conventional quantile regression gives only estimates that are 0.00 or
+1.00 which are artifacts of the rounding error of gestational age.
Certainly some caution in interpretation is warranted regarding the role
of random sampling error in the estimates. If we assume that the median
duration of gestation is 40 weeks (280 days), then the standard error
comparing equal-sized groups of females versus males \(n_0=n_1=1250\) at
the median is roughly equal to \(\sqrt{\xi ^2(1/n_1 + 1/n_2)}\), where
\(\xi^2 = \tau(1-\tau)/f(280 \textrm{ days})^2\). If the density at 280
days is 2\%, then the standard error is approximately equal to 1. Thus
the results for female gender are imprecise. However, imprecise results
from random sampling error does not justify poor accuracy of point
estimation.

\section{Concluding Remarks}\label{concluding-remarks}

A critique of the preceding analysis is that quantile regression can
only be used to analyze continuous outcomes, and it should not be
applied to discrete outcomes. This is described, for example, in Koenker
(Koenker 2005), Section 3.2, and Davino, Furno and Vistocc (Davino et
al. 2014), Section 1.1.1, which stipulate that the distribution of the
outcome \(Y\) should have a continuous probability density. This
requirement justifies the large-sample confidence intervals and p-values
for quantile regression. Furthermore, when quantile regression is
applied to large datasets with continuous outcomes, then implausible
results that we saw in Section 4 are unusual because the gaps (spacings)
between neighboring data points shrink to zero. But on the other hand,
the process of fitting a quantile regression model by minimizing the
pinball loss in equation Equation~\ref{eq-fqr-QRobjFunction} does not
actually require that the outcome \(Y\) is continuous. To put it
bluntly: Equation Equation~\ref{eq-fqr-QRobjFunction} doesn't care if
\(Y\) is continuous or not.

Moreover, the requirement of a continuous outcome in quantile regression
is limiting in practice. In many applications the outcome variable is
not truly continuous. In statistics, the word ``continuous'' is used in
a narrow mathematical sense for a scale in which any real number in its
range is a possible value. However, the word ``continuous'' in a
practical sense means a scale that is very fine for the application at
hand. In the PMS data example, the true gestational age of an infant is
a continuous quantity: It is the duration of human pregnancy from
conception until birth of the infant. But in practice, gestational age
is rarely measured with precision exceeding \(\pm 1\) day because the
exact timing of the start of conception is difficult to measure
precisely.

And yet, despite the discreteness of \(Y\) that arises from rounding
errors in the data, epidemiologists still wish to use quantile
regression for outcomes like gestational age. The method provides
information about the entire distribution of the response, and this
allows for a deeper understanding of how predictor variables influence
different parts of the outcome distribution, not just the average. The
benefits of quantile regression are clearly displayed in Table
Table~\ref{tbl-fqr-pmsbetasqr}: We see that albuminuria has the largest
impact at the lower end of the gestational age distribution. This
finding invites inquiry into understanding the reasons for this large
effect. For example, some infants in the PMS study may be uniquely
susceptible to maternal illness during pregnancy.

There are different strategies to fix the implausible results in
quantile regression that result from rounding, and this is an ongoing
area of research (e.g.~see (Alampi et al. 2025)). Section 5 uses
dithering (Machado and Silva 2005), which adds uniform noise to the
outcome variable \(Y\) to account for the discreteness. In fact,
dithering can be viewed as an example of multiple imputation for missing
data. In multiple imputation, we replace a single missing data point
with one of several plausible, simulated values based on observed data
using Monte Carlo simulation. The process can be repeated several times
to average over the imputation uncertainty. Dithering can also be
interpreted as an example of ordered categorical regression with a
latent variable and cutpoints. Under this interpretation, the observed
discrete outcome is viewed as a rounded or categorized version of an
underlying continuous latent variable, and the added uniform noise
approximates the uncertainty regarding the true value within each
interval defined by the cutpoints. A different strategy to account for
implausible results is to use Bayesian quantile regression (Yang et al.
2016). This approach uses an asymmetric Laplace likelihood, and model
fitting proceeds using Markov chain Monte Carlo instead of the usual
pinball loss minimization. The advantage of Bayesian quantile regression
is that it explores the entire parameter space instead of being
restricted to integer-valued coefficients, and the result is the full
posterior distribution for unknown parameters.

\section*{References}\label{references}
\addcontentsline{toc}{section}{References}

\protect\phantomsection\label{refs}
\begin{CSLReferences}{1}{1}
\bibitem[\citeproctext]{ref-alampi2025performance}
Alampi, Jessica D., Bruce P. Lanphear, and Lawrence C. McCandless. 2025.
{``Performance of Quantile Regression Methods with Discrete Outcomes: A
Simulation Study with Applications to Environmental Epidemiology.''}
\emph{Environmental Epidemiology} 9 (6): e432.
\url{https://doi.org/10.1097/EE9.0000000000000432}.

\bibitem[\citeproctext]{ref-NCDS_UKDS}
Centre for Longitudinal Studies. 2024. {``National Child Development
Study (NCDS), 1958--2024.''} UK Data Service.
\url{https://doi.org/10.5255/UKDA-Series-2000032}.

\bibitem[\citeproctext]{ref-davino2014quantile}
Davino, Cristina, Marilena Furno, and Domenico Vistocco. 2014.
\emph{Quantile Regression: Theory and Applications}. Wiley.

\bibitem[\citeproctext]{ref-Divon2002}
Divon, M. Y., M. Ferber, R. Nisell, and B. Westgren. 2002. {``Male
Gender Predisposes to Prolongation of Pregnancy.''} \emph{American
Journal of Obstetrics and Gynecology} 187 (4): 1081--83.
\url{https://doi.org/10.1067/mob.2002.127141}.

\bibitem[\citeproctext]{ref-goldenberg2008epidemiology}
Goldenberg, Robert L., Jennifer F. Culhane, Jay D. Iams, and Roberto
Romero. 2008. {``Epidemiology and Causes of Preterm Birth.''} \emph{The
Lancet} 371 (9606): 75--84.
\url{https://doi.org/10.1016/S0140-6736(08)60074-4}.

\bibitem[\citeproctext]{ref-koenker2005quantile}
Koenker, Roger. 2005. \emph{Quantile Regression}. Cambridge University
Press.

\bibitem[\citeproctext]{ref-machado2005quantiles}
Machado, José A. F., and J. M. C. Santos Silva. 2005. {``Quantiles for
Counts.''} \emph{Journal of the American Statistical Association} 100
(472): 1226--37. \url{https://doi.org/10.1198/016214505000000330}.

\bibitem[\citeproctext]{ref-strickland2009ambient}
{Strickland, Matthew J. et al.} 2009. {``Associations Between Ambient
Air Pollution and Birth Weight.''} \emph{Unknown}.

\bibitem[\citeproctext]{ref-Yang2016Posterior}
Yang, Yunwen, Huixia Judy Wang, and Xuming He. 2016. {``Posterior
Inference in Bayesian Quantile Regression with Asymmetric Laplace
Likelihood.''} \emph{International Statistical Review} 84 (3): 327--44.
\url{https://doi.org/10.1111/insr.12114}.

\end{CSLReferences}




\end{document}
